% ============================================
% Chapter 5: Software Testing and Analysis Report
% ============================================
\chapter{Software Testing and Analysis Report}
\label{ch:testing-report}

This chapter reports system validation from two sources: (1) baseline capability tests inherited from the 2025 summer research report, and (2) the current skill-framework-specific evaluation added in this thesis. Following this arrangement, Sections 5.1--5.5 present the inherited baseline tests, while Section 5.6 reports the skill framework test retained in this work.

\section{Experiment Setup (Inherited Baseline)}
\label{sec:test-env}

\subsection{Hardware and Software Environment}
Baseline experiments were conducted in a Python 3.8 environment with the RS-agent stack integrated with RSHub 2.0. The reported environment includes Xeon-class server resources, GPU acceleration, FastAPI/LangChain service stack, and FAISS-based retrieval support.

\subsection{Datasets and Test Cases}
The inherited test suite includes:
\begin{itemize}
	\item knowledge corpus and workflow documents for snow, soil, and vegetation scenarios
	\item a natural-language query set for intent classification and domain QA
	\item multi-scenario modeling tasks for RSHub submission and completion validation
\end{itemize}

\section{Intent Recognition Evaluation}
\label{sec:test-intent}

\subsection{Objective}
Evaluate whether the agent can correctly classify user requests into major interaction intents, including casual chat, knowledge QA, data modeling, and parameter inversion.

\subsection{Results}
According to the inherited report, overall intent-recognition accuracy reaches \textbf{86.2\%}. Performance by intent type is summarized in Table~\ref{tab:intent-eval}.

\begin{table}[htbp]
	\centering
	\caption{Intent recognition performance (inherited baseline)}
	\label{tab:intent-eval}
	\begin{tabular}{lll}
		\hline
		Intent Type & Accuracy & Avg. Response Time (s) \\
		\hline
		Casual chat (0) & 95.0\% & 1.26 \\
		Knowledge QA (1) & 100.0\% & 1.26 \\
		Data modeling (2) & 80.0\% & 1.26 \\
		Parameter inversion (3) & 70.0\% & 1.26 \\
		Overall average & 86.2\% & 1.26 \\
		\hline
	\end{tabular}
\end{table}

\section{Domain Knowledge QA Evaluation}
\label{sec:test-qa}

\subsection{Objective}
Assess QA quality under RAG-based knowledge retrieval for microwave remote sensing topics.

\subsection{Results}
The inherited report shows an overall QA accuracy of \textbf{80.2\%}. Detailed performance is shown in Table~\ref{tab:qa-eval}.

\begin{table}[htbp]
	\centering
	\caption{Knowledge QA performance by domain (inherited baseline)}
	\label{tab:qa-eval}
	\begin{tabular}{llll}
		\hline
		Knowledge Domain & Accuracy & Completeness & Relevance \\
		\hline
		Basic theory & 85.0\% & 82.1\% & 83.5\% \\
		Model principles & 81.2\% & 79.4\% & 80.3\% \\
		Parameter setting & 78.9\% & 76.7\% & 77.8\% \\
		Application scenarios & 76.8\% & 74.2\% & 75.5\% \\
		Result analysis & 79.1\% & 77.3\% & 78.2\% \\
		Overall average & 80.2\% & 77.9\% & 79.1\% \\
		\hline
	\end{tabular}
\end{table}

\section{RSHub Submission and Execution Evaluation}
\label{sec:test-submit}

\subsection{Objective}
Validate integration behavior of RS-agent with RSHub platform, including task submission, parameter parsing, and completion status.

\subsection{Results}
Across snow/soil/vegetation scenarios, the inherited report records \textbf{100\%} submission success, \textbf{90.0\%} average parameter parsing accuracy, and \textbf{80.6\%} average task completion rate.

\begin{table}[htbp]
	\centering
	\caption{RSHub task submission performance (inherited baseline)}
	\label{tab:submit-eval}
	\begin{tabular}{llll}
		\hline
		Scene Type & Submission Success & Parameter Parsing Accuracy & Task Completion Rate \\
		\hline
		Snow & 100\% & 89.9\% & 79.9\% \\
		Soil & 100\% & 90.0\% & 80.8\% \\
		Vegetation & 100\% & 90.0\% & 81.0\% \\
		Overall average & 100\% & 90.0\% & 80.6\% \\
		\hline
	\end{tabular}
\end{table}

\section{Modeling Capability and End-to-End Integration}
\label{sec:test-e2e}

\subsection{Modeling Capability Results}
For representative modeling cases across snow/soil/vegetation, the inherited report indicates overall scores around mid-70\% level, with stronger behavior on snow/soil cases and relatively weaker performance on complex vegetation scenarios.

\begin{table}[htbp]
	\centering
	\caption{Modeling capability evaluation (inherited baseline)}
	\label{tab:modeling-eval}
	\begin{tabular}{lllll}
		\hline
		Case & Parameter Rationality & Result Accuracy & Computational Efficiency & Visualization Quality \\
		\hline
		Snow modeling & 76.8\% & 74.2\% & 75.1\% & 78.9\% \\
		Soil modeling & 77.5\% & 75.8\% & 74.3\% & 81.2\% \\
		Vegetation modeling & 70.7\% & 71.9\% & 72.8\% & 68.8\% \\
		Overall average & 75.0\% & 74.0\% & 74.1\% & 76.3\% \\
		\hline
	\end{tabular}
\end{table}

\subsection{Integrated Workflow Results}
End-to-end workflow evaluation in the inherited report gives a task completion level of \textbf{91.2\%}, with stable multi-step execution in representative remote sensing workflows.

\begin{table}[htbp]
	\centering
	\caption{System integration performance (inherited baseline)}
	\label{tab:integration-eval}
	\begin{tabular}{lll}
		\hline
		Metric & Value & Grade \\
		\hline
		Task completion & 91.2\% & Excellent \\
		System stability & 88.5\% & Good \\
		Result reliability & 89.7\% & Excellent \\
		Execution efficiency & 85.4\% & Good \\
		\hline
	\end{tabular}
\end{table}

\section{Skill Framework and Token Optimization Test}
\label{sec:test-skill-token}

This section is retained as the thesis-specific evaluation of the current skill framework, and is intentionally placed at the end of this chapter.

\subsection{Objective}
Evaluate the effect of replacing monolithic prompts with tiered skill loading (layer-1 catalog + selective layer-2 full documents).

\subsection{Observed Results}
\begin{table}[H]
	\centering
	\caption{Skill architecture comparison (baseline vs tiered loading)}
	\label{tab:skill-token}
	{\footnotesize
		\setlength{\tabcolsep}{2.5pt}
		\renewcommand{\arraystretch}{1.1}
		\begin{tabular}{p{2.4cm}p{2.6cm}p{2.6cm}p{4.4cm}p{1.1cm}}
			\hline
			Scenario & Baseline (long prompt) & Tiered skill loading & Observation & Pass/Fail \\
			\hline
			Single-turn QA & Works & Works & No quality regression observed & Pass \\
			DMRT paper summary & Works & Works & No quality regression observed & Pass \\
			Task planning & Works & Works & No quality regression observed & Pass \\
			Multi-turn continuity & Context loss in some A-tests & Stable & Tiered mode improves context retention & Pass \\
			\hline
		\end{tabular}
	}
\end{table}
\noindent\textit{Measurement note: quantitative token/latency instrumentation is still limited in this run. The most visible difference is multi-turn context continuity under long interaction sequences.}

\section{Chapter Summary}

This chapter combines inherited baseline evaluation and current thesis-specific evaluation.
The inherited baseline (Sections 5.1--5.5) confirms practical capability in intent recognition, domain QA, RSHub task integration, modeling workflow execution, and end-to-end stability.
The retained thesis-specific test (Section 5.6) further shows that the skill framework improves context continuity while preserving response quality.
Overall, these results support the project-level claim that RSHub-agent is operationally viable and extensible for practical remote sensing workflows.
